\documentclass[a4paper,12pt]{article}

\usepackage[utf8]{inputenc}
\usepackage[T1]{fontenc}
\usepackage[italian]{babel}
\usepackage{geometry}
\usepackage{amsmath}
\usepackage{booktabs}
\usepackage{xcolor}
\usepackage{listings}
\usepackage{enumitem}

\geometry{margin=2.5cm}

\lstset{
    basicstyle=\ttfamily\small,
    backgroundcolor=\color{gray!10},
    frame=single,
    breaklines=true,
    columns=fullflexible
}

\title{Analisi dei parametri di stop di QALS}
\author{}
\date{}

\begin{document}

\maketitle

\section{Osservazione sui parametri \texttt{i\_max}, \texttt{N\_max} e \texttt{d\_min}}

Spulciando meglio i parametri di QALS, mi sono accorto che era già impostato un parametro relativo alle iterazioni di QALS: \texttt{i\_max}. 
Il parametro in questione rappresenta il numero massimo di iterazioni di QALS, cioè il limite superiore della ricerca. In altre parole, rappresenta quanto si lascia che QALS cerchi soluzioni.

Il comportamento di \texttt{i\_max} può essere riassunto nel seguente modo:

\begin{itemize}[leftmargin=1.5cm]
    \item \textbf{\texttt{i\_max} basso}: QALS è molto più veloce (a causa delle meno iterazioni) ma le soluzioni che si vanno a trovare 
                                          sono peggiori inquanto QALS non aggiorna abbastanza la tabu matrix e non riesce ad esplorare bene
                                          lo spazio delle soluzioni

    \item \textbf{\texttt{i\_max} alto}: QALS è chiaramente più lento ma a discapito della velocità riesce ad esplorare meglio lo spazio delle
                                         soluzioni e ad aggiornare in modo efficace la tabu matrix, qundi le soluzioni dovrebbero risultare migliori.
\end{itemize}

Il parametro \texttt{i\_max} è infatti uno dei parametri che definisce la condizione di stop di QALS:

\begin{lstlisting}[language=Python]
if ((i == i_max) or ((e + d >= N_max) and (d < d_min))):
    . . .
    break
\end{lstlisting}

Questa condizione permette di verificare se QALS ha raggiunto il numero massimo di iterazioni oppure 
se è entrato in una condizione di convergenza o stallo.

\section{Proposta modifica}

Alla luce di questa osservazione, quello che suggerisco è di non aumentare ulteriormente il dominio di \texttt{TIMES}. Al contrario, abbasserei \texttt{TIMES}, per esempio fissandolo al solo valore:

\[
\texttt{TIMES} \in \{10\}
\]

D'altro canto, aumenterei invece il dominio di \texttt{i\_max}, in modo da studiare meglio quanto la profondità della ricerca interna di QALS influenzi la qualità delle soluzioni trovate:

\[
\texttt{i\_max} \in \{10, 50, 100, 250, 500, 1000\}
\]

In questo modo, \texttt{TIMES} viene utilizzato solo per ottenere una piccola statistica sulle run indipendenti, mentre \texttt{i\_max} diventa il vero parametro con cui si controlla la profondità della ricerca di QALS.

\section{Altri parametri della condizione di stop}

Gli altri parametri di QALS che definiscono la condizione di stop sono \texttt{N\_max} e \texttt{d\_min}.

\begin{itemize}[leftmargin=1.5cm]
    \item \textbf{\texttt{N\_max}}: numero massimo di iterazioni se l'alg non migliora.
                                    Proporrei di impostare un dominio proporzionale a \texttt{i\_max}:
    \[
    \texttt{N\_max} \in \{5, 25, 50, 125, 250, 500\}
    \]

    \item \textbf{\texttt{d\_min}}: conta quante volte trovi una soluzione diversa ma peggiore della migliore corrente. 
                                    Per una prima prova, imposterei:
    \[
    \texttt{d\_min} = 0.7 \cdot \texttt{N\_max}
    \]
    ottenendo quindi il seguente dominio:
    \[
    \texttt{d\_min} \in \{4, 18, 35, 88, 175, 350\}
    \]
\end{itemize}

\section{Tabella dei domini proposti}

\begin{table}[h]
\centering
\begin{tabular}{ccc}
\toprule
\texttt{i\_max} & \texttt{N\_max} & \texttt{d\_min} \\
\midrule
10   & 5   & 4   \\
50   & 25  & 18  \\
100  & 50  & 35  \\
250  & 125 & 88  \\
500  & 250 & 175 \\
1000 & 500 & 350 \\
\bottomrule
\end{tabular}
\caption{Domini proposti per i parametri di stop di QALS.}
\end{table}

\section{Configurazione precedente}

I parametri fino ad ora erano settati nel seguente modo:

\begin{itemize}[leftmargin=1.5cm]
    \item \texttt{i\_max = 10}
    \item \texttt{N\_max = 50}
    \item \texttt{d\_min = 70}
\end{itemize}

Questa configurazione è sbilanciata rispetto a \texttt{i\_max}, perché \texttt{N\_max} 
e \texttt{d\_min} sono maggiori del numero massimo di iterazioni consentite. 
Di conseguenza, con \texttt{i\_max = 10}, l'arresto dell'algoritmo è dominato quasi 
esclusivamente dal limite massimo di iterazioni, senza sfruttare davvero la condizione 
di stop basata su \texttt{i\_max}, \texttt{N\_max} e \texttt{d\_min}.

\end{document}